\input{\string~/Documents/School/"UC Irvine"/ta_template2.0.tex}
\rh{2B}{11.10}
\chead{\today}
\graphicspath{ {\string~/Desktop} }
\usetikzlibrary{arrows}
%============ANSWER KEY TOGGLE===========
\newcommand{\ans}[1]{ \noindent {\color{red} \textbf{Answer:} #1}}
\renewcommand{\ans}[1]{\vspace{3in}}
%==========================================
\begin{document}
\begin{center}
\textbf{Taylor and Maclaurin Series}
\end{center}

\begin{thm}[Taylor Expansions]
	If $f$ has a power series representation at $a$, that is, there is an $R$ such that
	\[
		f(x) = \sum_{n = 0}^\infty c_n (x - a)^n \mte{ for } \lrabs{x - a} < R
	\]
	then the coefficients are given by:
	\[
			c_n := \frac{f^{(n)}(a)}{n!} 
	\]
	The series representation:
	\[
		f(x) = \sum_{n = 0}^\infty \frac{f^{(n)}(a)}{n!} (x-a)^n
	\]
	is called the \emph{Taylor series of the function $f$ at $a$}. When $a=0$, we call this the Maclaurin series.
\end{thm}
	


\begin{thm}
[Remainder Theorem]
	If $f(x)=T_n(x)+R_n(x)$, where $T_n$ is the $n$ th-degree Taylor polynomial of $f$ at $a$ and
$$
\lim _{n \rightarrow \infty} R_n(x)=0
$$
for $|x-a|<R$, then $f$ is equal to the sum of its Taylor series on the interval $|x-a|<R$
\end{thm}

\begin{thm}[Taylor's Inequality]
	If $\lrabs{f^{n+1}(x)} \leq M$ for $\lrabs{x-a} \leq d$, then the remainder $R_n(x)$ of the Taylor series satisfies the inequality:
	\[
		\lrabs{R_n(x)} \leq \frac{M}{(n+1)!} \lrabs{x-a}^{n+1}
	\]
	for $\lrabs{x-a} \leq d$
\end{thm}

\begin{thm}
		The Taylor expansions important functions are below:
		\[
\begin{array}{|c|c|c|}
\hline
\textbf{Function} & \textbf{Taylor Series} & \textbf{Expansion} \\ \hline
\frac{1}{1-x} & \sum_{n=0}^{\infty} x^n & 1+x+x^2+x^3+\cdots \\

e^x & \sum_{n = 0}^\infty \frac{x^n}{n!} & 1 + x + \frac{x^2}{2} + \frac{x^3}{3!} + \cdots \\ 
\sin (x) & \sum_{n=0}^{\infty}(-1)^n \frac{x^{2 n+1}}{(2 n+1) !} & x-\frac{x^3}{3 !}+\frac{x^5}{5 !}-\frac{x^7}{7 !}+\cdots \\
\cos(x) & \sum_{n=0}^{\infty}(-1)^n \frac{x^{2 n}}{(2 n) !} & 1-\frac{x^2}{2 !}+\frac{x^4}{4 !}-\frac{x^6}{6 !}+\cdots \\ 
\tan { }^{-1} x & \sum_{n=0}^{\infty}(-1)^n \frac{x^{2 n+1}}{2 n+1} & x-\frac{x^3}{3}+\frac{x^5}{5}-\frac{x^7}{7}+\cdots \\
 \ln (1+x) & \sum_{n=1}^{\infty}(-1)^{n-1} \frac{x^n}{n} & x-\frac{x^2}{2}+\frac{x^3}{3}-\frac{x^4}{4}+\cdots \\
\hline
\end{array}
		\]
\end{thm}

\begin{exx}
	Provide a Taylor series for the function $f(x) = xe^x$ centered at $a = 0$
\end{exx}
\ans{Take the series from the table above and multiply by $x$:
\[
\sum_{n = 0}^\infty \frac{x^{n+1}}{n!}
\]
}

\problembox{Find the first four terms of the Taylor series centered at $a = 2$ of the function:
\[
f(x) = \frac{1}{1 + x} 
\]}
\ans{This is just computing a lot of derivatives.}

\problembox{Find the Taylor Series for $f(x)=\ln (3+4 x)$ about $x=0$.}
\ans{\href{https://tutorial.math.lamar.edu/Solutions/CalcII/TaylorSeries/Prob4.aspx}{Link} }

\begin{aun}
	Here's another \href{https://math.libretexts.org/Courses/Monroe_Community_College/MTH_211_Calculus_II/Chapter_Ch10%3A_Power_Series/10.3%3A_Taylor_and_Maclaurin_Series/10.3E%3A_Exercises_for_Taylor_Polynomials_and_Taylor_Series}{helpful link}. 
\end{aun}



\end{document}